% ৪.৭ HTML-এর প্রাথমিক ধারণা
\section{HTML-এর প্রাথমিক ধারণা}

\begin{learningbox}
এই টপিক শেষে শিক্ষার্থী HTML-এর সংজ্ঞা, পূর্ণরূপ, ইতিহাস, HTML5, file extension, tag/element/attribute-এর মৌলিক ধারণা ব্যাখ্যা করতে পারবে।
\end{learningbox}

\subsection{ভূমিকা}
ইন্টারনেটে দেখা web page-এর structure ও content presentation তৈরি করতে HTML ব্যবহার করা হয়। HTML tag browser-কে বলে কোন অংশ heading, কোন অংশ paragraph, কোন অংশ image বা link হবে।

\subsection{HTML-এর পূর্ণরূপ}
\begin{definitionbox}{HTML}
HTML = HyperText Markup Language
\end{definitionbox}

\begin{keypointbox}{HyperText, Markup, Language}
\begin{itemize}
\item HyperText: এমন text যাতে link থাকে এবং এক page থেকে অন্য page-এ যাওয়া যায়।
\item Markup: tag ব্যবহার করে content চিহ্নিত ও format করা।
\item Language: web page তৈরির জন্য ব্যবহৃত markup system।
\end{itemize}
\end{keypointbox}

\subsection{HTML কি programming language?}
\begin{mistakebox}{উল্লেখ্য}
HTML programming language নয়; এটি একটি markup language। HTML কোন logic বা calculation করে না; presentation ও behaviour-এর জন্য CSS ও JavaScript ব্যবহৃত হয়।
\end{mistakebox}

\subsection{HTML দিয়ে কী করা যায় ও না করা যায়}
\begin{keypointbox}{কি করা যায়}
\begin{itemize}
\item Heading, paragraph, image, hyperlink, list, table, form তৈরি করা।
\item Page-এর basic structure নির্ধারণ করা।
\end{itemize}
\end{keypointbox}

\begin{keypointbox}{কি করা যায় না (এককভাবে)}
HTML দিয়ে single-handedly dynamic features, database connection, authentication বা complex logic করা যায় না; এসবের জন্য server-side language, database ও client-side scripting প্রয়োজন।
\end{keypointbox}

\subsection{সংক্ষিপ্ত ইতিহাস}
Tim Berners-Lee 1990-এর দিকে HTML-এর প্রাথমিক ভাবনা উপস্থাপন করেন। পরবর্তীতে HTML-এর বিভিন্ন সংস্করণ আসে; আধুনিক সময়ের কাজগুলো HTML5 ও HTML Living Standard-এ রূপ নেয়।

\subsection{HTML5}
HTML5 হল HTML-এর আধুনিক মান/সংস্করণ, যা multimedia (audio/video), semantic elements (header, nav, main, article, footer), improved form controls ইত্যাদি সমর্থন করে।

\subsection{File extension}
HTML file সাধারণত \texttt{.html} বা \texttt{.htm} extension দিয়ে সংরক্ষণ করা হয়। প্রাথমিক পেজের নাম সাধারণত \texttt{index.html} হয়।

\subsection{HTML Tag, Element ও Attribute (সংক্ষেপ)}
\begin{description}
\item[Tag] Angle bracket \(< >\) এর মধ্যে লেখা keyword, যেমন \texttt{<p>}।
\item[Element] Opening tag, content ও closing tag মিলে element গঠন করে, উদাহরণ: \texttt{<p>This is a paragraph.</p>}।
\item[Attribute] Tag-এর অতিরিক্ত তথ্য, যেমন \texttt{src}, \texttt{alt}, \texttt{width} ইত্যাদি: \texttt{<img src="photo.jpg" alt="Photo">}।
\end{description}

\subsection{Simple HTML example}
\begin{examplebox}{Code}
\begin{verbatim}
<!DOCTYPE html>
<html>
  <head>
    <title>My First Web Page</title>
  </head>
  <body>
    <h1>Welcome to HSC ICT</h1>
    <p>This is my first HTML web page.</p>
  </body>
</html>
\end{verbatim}
\end{examplebox}

\subsection{Common mistakes}
\begin{itemize}
\item Closing tag ভুলে যাওয়া
\item Tag-এর জন্য angle bracket নয় অন্য bracket ব্যবহার করা
\item File extension ভুল (\texttt{.txt} হিসেবে save করা)
\item Attribute value-তে quotation না দেওয়া
\end{itemize}

\begin{practicebox}
\begin{enumerate}
\item HTML-এর পূর্ণরূপ লেখ।
\item একটি ছোট HTML skeleton লিখে browser-এ খুলে দেখো।
\item HTML, CSS ও JavaScript-এর ভূমিকা সংক্ষেপে লেখো।
\end{enumerate}
\end{practicebox}
